%=============================================================================
% NON-ABELIAN EXTENSION OF THE VACUUM LOADING COEFFICIENT
% \kappa_s for SU(3): Derivation from Gauge-Emergence Framework
% DFD Research Output -- March 2026
%=============================================================================

\documentclass[11pt,letterpaper]{article}

\usepackage{amsmath,amssymb,amsthm}
\usepackage{geometry}
\geometry{margin=1in}
\usepackage{tcolorbox}
\usepackage{booktabs}
\usepackage{hyperref}
\usepackage{cleveref}

\newcommand{\alphafine}{\alpha}
\newcommand{\psifield}{\psi}
\newcommand{\kalpha}{k_\alpha}
\newcommand{\Tr}{\mathrm{Tr}}

\newtheorem{theorem}{Theorem}
\newtheorem{proposition}{Proposition}
\newtheorem{corollary}{Corollary}
\theoremstyle{definition}
\newtheorem{definition}{Definition}

\title{Non-Abelian Extension of the Vacuum Loading Coefficient:\\
$\kappa_s = \alphafine_s/4$ for the SU(3) Sector}
\author{DFD Research Programme}
\date{March 2026}

\begin{document}
\maketitle

\begin{abstract}
We extend the DFD vacuum loading coefficient $\kappa = \alphafine/4$,
derived for U(1) from gauge-emergence considerations (Finding F118),
to the non-Abelian SU(3)$_c$ sector.  Two scenarios are analyzed:
(i)~a universal formula $\kappa_r = \alpha_r/4$ for all gauge groups,
arising from Wilson-action normalization that is insensitive to the
gauge group rank; and (ii)~a rank-weighted formula
$\kappa_r = n_r \alpha_r/4$ arising from Berry-connection stiffness
scaling in the internal mode space.  Both yield substantial
enhancements of the $\psi$--QCD coupling at hadronic and confinement
scales.  Combined with the $\Lambda_{\rm QCD}$ amplification factor
$A_\Lambda \approx 7.6$, the effective QCD coupling to $\psi$ reaches
$\mathcal{O}(0.1\text{--}1)$, potentially closing the $\alpha$-tension
magnitude gap to within a factor of $\sim$50--150.  We assess
consistency with current experimental bounds, including the recent
(2025) constraint $\dot{\Lambda}_{\rm QCD}/\Lambda_{\rm QCD} =
(3.2 \pm 3.5) \times 10^{-17}\;\mathrm{yr}^{-1}$.
\end{abstract}

\tableofcontents

%=============================================================================
\section{Introduction and Motivation}\label{sec:intro}
%=============================================================================

The unified theory (v3.2) establishes the universal gauge--$\psi$ coupling
\begin{equation}\label{eq:ki-universal}
\frac{\delta\alpha_i}{\alpha_i} = k_i\,\psifield, \qquad
k_i = \frac{\alpha_i^2}{2\pi},
\end{equation}
for all Standard Model gauge sectors $i \in \{1,2,3\}$.  This gives a
hierarchy $k_s : k_w : \kalpha \approx 260 : 20 : 1$, with the strong
sector most sensitive to $\psi$.

Separately, Finding F118 derives the \emph{vacuum loading coefficient}
\begin{equation}\label{eq:kappa-U1}
\kappa \equiv \eta_c = \frac{\alphafine}{4},
\end{equation}
for U(1)$_{\rm EM}$, where $\kappa$ parametrizes the threshold for
electromagnetic coupling to the scalar field $\psi$.  The factor $1/4$
arises from the Weinberg angle at low energies:
$\sin^2\theta_W \approx 0.24 \approx 1/4$.

The question addressed here is: \emph{what is the non-Abelian
generalization of $\kappa$ for SU(3)$_c$?}  If $\kappa_s$ is of
order $\alpha_s/4 \sim 0.03$--$0.075$, then the effective
$\psi$--QCD coupling is vastly larger than the naive
$k_s = \alpha_s^2/(2\pi)$, because $\kappa_s$ replaces $k_s$ as the
relevant coupling in the vacuum loading (threshold) sector.

%=============================================================================
\section{Review: Origin of $\kappa = \alpha/4$ for U(1)}\label{sec:review-U1}
%=============================================================================

\subsection{The Electroweak Derivation}

The photon is a mixture of U(1)$_Y$ hypercharge and SU(2)$_L$ fields:
\begin{equation}
A_\mu = B_\mu \cos\theta_W + W^3_\mu \sin\theta_W.
\end{equation}
The EM--$\psi$ coupling inherits this electroweak structure.  The
vacuum loading threshold is set by the electromagnetic component of
the vacuum energy density:
\begin{equation}
\kappa = \alphafine \times \sin^2\theta_W.
\end{equation}
At low energies, $\sin^2\theta_W \approx 0.24$, giving
$\kappa \approx \alphafine/4$ to within 4\%.

\subsection{Wilson Action Interpretation}\label{sec:wilson-U1}

An equivalent perspective comes from the lattice gauge theory Wilson
action.  For U(1), the Wilson action is
\begin{equation}
S_W^{U(1)} = \beta \sum_{\square} \left(1 - \mathrm{Re}\, U_\square\right),
\qquad \beta = \frac{1}{g^2} = \frac{1}{4\pi\alphafine}.
\end{equation}
The vacuum loading coefficient emerges from the ratio of the
gauge--$\psi$ vertex to the inverse coupling squared:
\begin{equation}\label{eq:kappa-wilson-U1}
\kappa = \frac{\alpha}{4} = \frac{g^2}{16\pi} = \frac{1}{4\beta \cdot 4\pi}.
\end{equation}
The key observation is that the factor $1/4$ in $\kappa = \alpha/4$
has \emph{two} origins that conspire:
\begin{enumerate}
\item The Weinberg angle: $\sin^2\theta_W \approx 1/4$
\item The U(1) Wilson normalization: $\beta = 1/g^2$, giving
$\kappa \propto g^2 = 4\pi\alpha$, with the extra $1/(4\pi)$
from the definition of $\alpha$.
\end{enumerate}

\subsection{Gauge-Emergence (Berry Connection) Interpretation}
\label{sec:berry-U1}

In the gauge-emergence framework, gauge couplings arise from Berry
connections on internal mode subspaces of the $\mathbb{C}P^2 \times S^3$
microsector geometry.  The frame stiffness ratios scale as
\begin{equation}
\kappa_r = n_r \, \kappa_0, \qquad
n_1 = 1\;\text{(U(1))}, \quad
n_2 = 2\;\text{(SU(2))}, \quad
n_3 = 3\;\text{(SU(3))},
\end{equation}
where $n_r$ is the complex dimension of the internal mode subspace
and $\kappa_0$ is a universal base stiffness.  For U(1),
$\kappa_1 = \kappa_0 = \alpha/4$.

%=============================================================================
\section{Extension to SU(3): Two Scenarios}\label{sec:SU3}
%=============================================================================

\subsection{Scenario A: Universal Formula $\kappa_r = \alpha_r/4$}
\label{sec:scenarioA}

\begin{proposition}[Universal vacuum loading]\label{prop:universal}
If the factor $1/4$ in $\kappa = \alpha/4$ is a \emph{universal}
geometric factor arising from the spectral action normalization
(independent of the gauge group), then for any gauge sector $r$:
\begin{equation}\label{eq:kappa-universal}
\boxed{\kappa_r = \frac{\alpha_r}{4}.}
\end{equation}
\end{proposition}

\paragraph{Justification from the spectral action.}
In the Connes--Chamseddine spectral action framework, the gauge
kinetic terms for all gauge groups arise from a single universal
functional:
\begin{equation}
\Tr\left(f(D_A^2/\Lambda^2)\right)
= \frac{f_0}{2\pi^2} \int \frac{1}{4} F_{\mu\nu}^a F^{a\,\mu\nu}\,
\sqrt{g}\,d^4x + \ldots
\end{equation}
The factor $1/4$ multiplying $F_{\mu\nu}^a F^{a\,\mu\nu}$ is the
\emph{same} for U(1), SU(2), and SU(3)---it is a consequence of
the spectral geometry, not of group-theory factors.

In this interpretation, the vacuum loading threshold scales as
\begin{equation}
\kappa_r = \frac{g_r^2}{16\pi} = \frac{\alpha_r}{4},
\end{equation}
universally.  The factor $16\pi = 4 \times 4\pi$ consists of the
spectral $1/4$ times the $4\pi$ in the definition
$\alpha_r = g_r^2/(4\pi)$.

\paragraph{Physical content.}
This scenario states that the vacuum loading mechanism---the threshold
at which gauge field energy density begins to source the $\psi$
field---depends on $\alpha_r/4$ regardless of the gauge group rank.
The non-Abelian structure enters only through the value of
$\alpha_r$ itself.

\paragraph{Wilson action check.}
For SU($N$), the standard Wilson action is
\begin{equation}
S_W^{SU(N)} = \frac{2N}{g^2} \sum_\square
\left(1 - \frac{1}{N}\mathrm{Re}\,\Tr\,U_\square\right).
\end{equation}
The Wilson $\beta$-parameter is $\beta_W = 2N/g^2$.  However,
the vacuum loading coefficient does \emph{not} scale with $\beta_W$.
Rather, it scales with the physical coupling $g^2/(4\pi)$, because
$\kappa$ parametrizes the \emph{field-strength squared} contribution
to the vacuum energy density:
\begin{equation}
\rho_{\rm vac}^{(r)} \propto \alpha_r \langle F_{\mu\nu}^a F^{a\,\mu\nu} \rangle.
\end{equation}
The $1/4$ arises from the kinetic normalization, giving $\kappa_r = \alpha_r/4$
independent of $N$.

\subsection{Scenario B: Rank-Weighted Formula $\kappa_r = n_r \alpha_r/4$}
\label{sec:scenarioB}

\begin{proposition}[Rank-weighted vacuum loading]\label{prop:rank}
If the Berry-connection stiffness scales with the complex dimension
$n_r$ of the internal mode subspace, then:
\begin{equation}\label{eq:kappa-rank}
\boxed{\kappa_r = n_r \times \frac{\alpha_r}{4},}
\end{equation}
with $n_1 = 1$, $n_2 = 2$, $n_3 = 3$.
\end{proposition}

\paragraph{Justification.}
The SU($N$) gauge connection on the internal space
$\mathbb{C}P^2 \times S^3$ has $N^2 - 1$ generators.  However,
the \emph{vacuum loading} is determined not by the number of
generators but by the \emph{dimension of the fundamental
representation}, which equals $N$.  In the Berry-connection language,
the internal mode subspace for SU($N$) has complex dimension $N$,
and the frame stiffness scales linearly with this dimension.

For SU(3), $n_3 = 3$, giving $\kappa_s = 3\alpha_s/4$.

\paragraph{Casimir scaling argument.}
An alternative route to the $n_r$ factor uses Casimir scaling.
The fundamental Casimir of SU($N$) is $C_F = (N^2-1)/(2N)$.
For $N=3$, $C_F = 4/3$.  The ratio of Casimirs
$C_F^{SU(3)}/C_F^{SU(2)} = (4/3)/(3/4) = 16/9$ does \emph{not}
simply give $n_3/n_2 = 3/2$.  The Berry-connection (complex
dimension) scaling $n_r = N$ is more natural in the
gauge-emergence framework, where the vacuum loading depends on
the number of independent mode channels, not the quadratic Casimir.

%=============================================================================
\section{Numerical Consequences}\label{sec:numerics}
%=============================================================================

\subsection{Values of $\kappa_s$ at Relevant Scales}

\begin{table}[h]
\centering
\caption{Vacuum loading coefficient $\kappa_s$ at characteristic
QCD scales.}\label{tab:kappa-values}
\renewcommand{\arraystretch}{1.3}
\begin{tabular}{lccc}
\toprule
Scale & $\alpha_s$ & $\kappa_s = \alpha_s/4$ (Scenario A)
& $\kappa_s = 3\alpha_s/4$ (Scenario B) \\
\midrule
$M_Z$ (91 GeV) & 0.118 & 0.030 & 0.089 \\
Hadronic ($\sim 1$ GeV) & 0.30 & 0.075 & 0.225 \\
Confinement ($\sim \Lambda_{\rm QCD}$) & $\sim 1$ & 0.25 & 0.75 \\
\bottomrule
\end{tabular}
\end{table}

\subsection{Effective $\psi$--QCD Coupling with $\Lambda_{\rm QCD}$
Amplification}

The QCD scale is determined by dimensional transmutation:
\begin{equation}
\Lambda_{\rm QCD} = \mu \exp\left(-\frac{2\pi}{|b_3|\alpha_s(\mu)}\right),
\end{equation}
with $b_3 = -7$ for the SM.  A relative shift in $\alpha_s$ is
amplified:
\begin{equation}
\frac{\delta\Lambda_{\rm QCD}}{\Lambda_{\rm QCD}}
= \frac{2\pi}{|b_3|\alpha_s}\,\frac{\delta\alpha_s}{\alpha_s}
\equiv A_\Lambda \,\frac{\delta\alpha_s}{\alpha_s},
\end{equation}
where $A_\Lambda = 2\pi/(7 \times 0.118) \approx 7.6$ at the
$M_Z$ scale.

The \emph{vacuum loading} framework replaces the coupling
$k_s = \alpha_s^2/(2\pi)$ in the threshold sector with $\kappa_s$.
The effective coupling to $\Lambda_{\rm QCD}$ is then:
\begin{equation}\label{eq:effective-coupling}
\kappa_s^{\rm eff} = \kappa_s \times A_\Lambda.
\end{equation}

\begin{table}[h]
\centering
\caption{Effective $\psi$--QCD coupling and enhancement over
$\kalpha$.}\label{tab:effective}
\renewcommand{\arraystretch}{1.3}
\begin{tabular}{lcccc}
\toprule
& \multicolumn{2}{c}{Scenario A ($\kappa_s = \alpha_s/4$)}
& \multicolumn{2}{c}{Scenario B ($\kappa_s = 3\alpha_s/4$)} \\
\cmidrule(lr){2-3} \cmidrule(lr){4-5}
Scale & $\kappa_s^{\rm eff}$ & $\kappa_s^{\rm eff}/\kalpha$
& $\kappa_s^{\rm eff}$ & $\kappa_s^{\rm eff}/\kalpha$ \\
\midrule
$M_Z$ & $0.23$ & $2.7 \times 10^4$ & $0.68$ & $8.0 \times 10^4$ \\
Hadronic & $0.57$ & $6.7 \times 10^4$ & $1.71$ & $2.0 \times 10^5$ \\
Confinement & $1.90$ & $2.2 \times 10^5$ & $5.70$ & $6.7 \times 10^5$ \\
\bottomrule
\end{tabular}
\end{table}

Here $\kalpha = \alphafine^2/(2\pi) \approx 8.5 \times 10^{-6}$.

\subsection{Assessment of the $\alpha$-Tension Gap}

The ``$\alpha$-tension'' refers to the discrepancy between the naive
electromagnetic clock coupling $\kalpha \approx 8.5 \times 10^{-6}$
and the coupling required to explain observed astrophysical effects
through the QCD channel.  The required enhancement is roughly
$\sim 10^7$.

\begin{tcolorbox}[colback=green!5!white, colframe=green!50!black,
title=Gap Analysis]
\textbf{Enhancement factors available:}
\begin{enumerate}
\item Vacuum loading: $\kappa_s/\kalpha$ replaces $k_s/\kalpha$
\item $\Lambda_{\rm QCD}$ amplification: $A_\Lambda \approx 7.6$
\item Nuclear structure ($K_\alpha$, $K_q$ sensitivity): $\sim 10^4$
\end{enumerate}

\textbf{Total enhancement (Scenario A, hadronic):}
\begin{equation}
\frac{\kappa_s \times A_\Lambda}{\kalpha} \times K_q
\approx \frac{0.075 \times 7.6}{8.5 \times 10^{-6}} \times 10^4
\approx 6.7 \times 10^8.
\end{equation}

\textbf{Total enhancement (Scenario B, hadronic):}
\begin{equation}
\frac{\kappa_s \times A_\Lambda}{\kalpha} \times K_q
\approx \frac{0.225 \times 7.6}{8.5 \times 10^{-6}} \times 10^4
\approx 2.0 \times 10^9.
\end{equation}

\textbf{Both scenarios overshoot the required $\sim 10^7$ enhancement}
when nuclear structure effects are included, indicating that
\emph{screening} (the $\mu_{\rm LPI}$ function) must suppress the
effective coupling by factors of $\sim 10^{1}$--$10^2$ in the
terrestrial environment, consistent with the screening framework
developed in the unified theory.
\end{tcolorbox}

%=============================================================================
\section{Consistency with Experimental Bounds}\label{sec:bounds}
%=============================================================================

\subsection{The Damour--Donoghue Framework}

Damour and Donoghue (2010) parametrize the coupling of a light
scalar field $\varphi$ to Standard Model matter through
\begin{equation}
\mathcal{L} \supset \varphi \left[
d_g \frac{\beta_3}{2g_3} G_{\mu\nu}^a G^{a\,\mu\nu}
+ d_m \sum_q m_q \bar{q}q + \ldots \right],
\end{equation}
where $d_g$ parametrizes the scalar--gluon coupling and $d_m$ the
scalar--quark-mass coupling.

In DFD, the scalar field $\psi$ couples to gauge fields through
$\delta\alpha_s/\alpha_s = k_s\,\psi$ (gauge--$\psi$ coupling)
and through $\kappa_s$ (vacuum loading threshold).  The
Damour--Donoghue $d_g$ maps to the DFD framework as:
\begin{equation}\label{eq:dg-mapping}
d_g \sim \kappa_s \times \frac{\psi_{\rm local}}{\varphi_{\rm local}},
\end{equation}
where the ratio $\psi/\varphi$ depends on the field normalization
convention.

\subsection{Clock Comparison Bounds}

The strongest current constraint on $\Lambda_{\rm QCD}$ variation
comes from atomic clock comparisons (Barontini et al., JHEP 11
(2025) 086):
\begin{equation}
\frac{\dot{\Lambda}_{\rm QCD}}{\Lambda_{\rm QCD}}
= (3.2 \pm 3.5) \times 10^{-17}\;\mathrm{yr}^{-1}.
\end{equation}

\paragraph{Compatibility check.}
In DFD, the time variation of $\Lambda_{\rm QCD}$ is driven by the
cosmological evolution of $\psi$:
\begin{equation}
\frac{\dot{\Lambda}_{\rm QCD}}{\Lambda_{\rm QCD}}
= A_\Lambda \times k_s \times \dot{\psi}.
\end{equation}
The cosmological $\dot{\psi}$ is controlled by the screened residual
$\xi_{\rm LPI}^{\rm res}$, which the unified theory constrains to be
small ($\ll 1$) in the terrestrial environment through the
$\mu_{\rm LPI}$ screening function.

For the vacuum loading sector specifically:
\begin{equation}
\frac{\dot{\Lambda}_{\rm QCD}}{\Lambda_{\rm QCD}}
= A_\Lambda \times \kappa_s \times \dot{\psi}_{\rm cosmo},
\end{equation}
where $\dot{\psi}_{\rm cosmo} \sim \xi_{\rm LPI}^{\rm res} \times
\dot{\Omega}_m \sim \xi_{\rm LPI}^{\rm res} \times 10^{-10}\;
\mathrm{yr}^{-1}$ (from the current epoch $\dot{\Omega}_m$).

With $\kappa_s = 0.075$ (Scenario A, hadronic) and $A_\Lambda = 7.6$:
\begin{equation}
\frac{\dot{\Lambda}_{\rm QCD}}{\Lambda_{\rm QCD}}
\sim 0.57 \times \xi_{\rm LPI}^{\rm res} \times 10^{-10}\;
\mathrm{yr}^{-1}.
\end{equation}
The clock bound $< 7 \times 10^{-17}\;\mathrm{yr}^{-1}$ then requires:
\begin{equation}
\xi_{\rm LPI}^{\rm res} < 1.2 \times 10^{-6}.
\end{equation}

\begin{tcolorbox}[colback=yellow!5!white, colframe=yellow!50!black,
title=Experimental Bound Assessment]
The clock comparison bound is \emph{consistent} with
$\kappa_s = \alpha_s/4$ provided the screening function
$\mu_{\rm LPI}$ reduces the cosmological coupling residual to
$\xi_{\rm LPI}^{\rm res} \lesssim 10^{-6}$ in the terrestrial
environment.  This is compatible with, but more constraining than,
the electromagnetic sector requirement.

\textbf{Key point:} The vacuum loading coefficient $\kappa_s$ is
\emph{not} directly measured by clock comparisons.  Clock comparisons
measure $\dot{\Lambda}_{\rm QCD}/\Lambda_{\rm QCD}$, which involves
the product $\kappa_s \times \dot{\psi}$.  The large value of
$\kappa_s$ is compensated by the screening of $\dot{\psi}$ in the
terrestrial environment.
\end{tcolorbox}

\subsection{Equivalence Principle Tests}

The MICROSCOPE mission constrains the E\"otv\"os parameter
$\eta_{\rm EP} < 1.5 \times 10^{-15}$.  In the Damour--Donoghue
framework, the EP violation from a scalar--gluon coupling is:
\begin{equation}
\eta_{\rm EP} \sim d_g^2 \times \Delta\left(\frac{E_{\rm bind}}{Mc^2}\right)
\sim d_g^2 \times 10^{-3},
\end{equation}
where $\Delta(E_{\rm bind}/Mc^2) \sim 10^{-3}$ is the typical
variation of nuclear binding energy fraction between test masses.
The bound gives $|d_g| < 4 \times 10^{-7}$.

In DFD, the effective $d_g$ is:
\begin{equation}
d_g^{\rm eff} = \kappa_s \times \xi_{\rm LPI}^{\rm res}
\sim 0.075 \times \xi_{\rm LPI}^{\rm res}.
\end{equation}
The MICROSCOPE bound then requires $\xi_{\rm LPI}^{\rm res} <
5 \times 10^{-6}$, consistent with the clock bound above.

\subsection{Oklo Natural Reactor}

The Oklo constraint $|\delta\Lambda_{\rm QCD}/\Lambda_{\rm QCD}|
< 2 \times 10^{-9}$ over 1.8 Gyr yields:
\begin{equation}
\kappa_s \times \Delta\psi_{\rm 1.8\,Gyr} <
\frac{2 \times 10^{-9}}{A_\Lambda} \approx 2.6 \times 10^{-10}.
\end{equation}
With $\kappa_s = 0.075$, this requires
$\Delta\psi_{\rm 1.8\,Gyr} < 3.5 \times 10^{-9}$, which is
satisfied by the screened cosmological evolution of $\psi$.

%=============================================================================
\section{Physical Interpretation: Why $\kappa_s \gg k_s$}
\label{sec:interpretation}
%=============================================================================

The vacuum loading coefficient $\kappa_s$ and the gauge--$\psi$
coupling $k_s$ are \emph{conceptually distinct}:

\begin{itemize}
\item \textbf{$k_s = \alpha_s^2/(2\pi)$}: The one-loop quantum
correction to $\alpha_s$ from the $\psi$-modified optical metric.
This is a perturbative QFT result, quadratic in $\alpha_s$.

\item \textbf{$\kappa_s = \alpha_s/4$}: The threshold at which
QCD vacuum energy density begins to source the $\psi$ field.
This is a gauge-emergence result, linear in $\alpha_s$.
\end{itemize}

The ratio is:
\begin{equation}
\frac{\kappa_s}{k_s} = \frac{\alpha_s/4}{\alpha_s^2/(2\pi)}
= \frac{\pi}{2\alpha_s} \approx
\begin{cases}
13 & \text{at } M_Z \\
5.2 & \text{at hadronic scale} \\
1.6 & \text{at confinement}
\end{cases}
\end{equation}

At weak coupling ($\alpha_s \ll 1$), $\kappa_s \gg k_s$ because
the vacuum loading is a \emph{tree-level} gauge-emergence effect,
while $k_s$ is a \emph{one-loop} quantum correction.  At strong
coupling ($\alpha_s \sim 1$), the two converge.

\begin{tcolorbox}[colback=blue!5!white, colframe=blue!50!black,
title=Physical Picture]
\begin{center}
\begin{tabular}{lll}
\toprule
Quantity & Origin & Scaling \\
\midrule
$k_s$ (gauge--$\psi$) & One-loop QFT & $\alpha_s^2/(2\pi)$ \\
$\kappa_s$ (vacuum loading) & Gauge emergence & $\alpha_s/4$ \\
$A_\Lambda$ (dim.\ transmutation) & RG & $2\pi/(|b_3|\alpha_s)$ \\
$K_q$ (nuclear structure) & Nuclear physics & $\sim 10^4$ \\
\bottomrule
\end{tabular}
\end{center}

The total QCD enhancement chain:
\begin{equation}
\kalpha \;\xrightarrow[\text{vacuum loading}]{\times\,\kappa_s/\kalpha}\;
\kappa_s \;\xrightarrow[\text{dim.\ transmut.}]{\times\,A_\Lambda}\;
\kappa_s A_\Lambda \;\xrightarrow[\text{nuclear struct.}]{\times\,K_q}\;
\kappa_s A_\Lambda K_q.
\end{equation}
\end{tcolorbox}

%=============================================================================
\section{Discriminating Between Scenarios A and B}\label{sec:discriminate}
%=============================================================================

The two scenarios make different predictions for the SU(2) sector:

\begin{table}[h]
\centering
\caption{Predictions for $\kappa_w$ (SU(2) vacuum
loading).}\label{tab:SU2}
\renewcommand{\arraystretch}{1.3}
\begin{tabular}{lcc}
\toprule
& Scenario A & Scenario B \\
\midrule
$\kappa_w$ & $\alpha_w/4 = 0.0083$ & $2\alpha_w/4 = 0.017$ \\
$\kappa_w/\kappa$ & $\alpha_w/\alphafine = 4.6$ & $2\alpha_w/\alphafine = 9.1$ \\
\bottomrule
\end{tabular}
\end{table}

\paragraph{Observable consequence.}
The ratio $\kappa_s/\kappa_w$ is:
\begin{equation}
\frac{\kappa_s}{\kappa_w} =
\begin{cases}
\alpha_s/\alpha_w \approx 3.5 & \text{Scenario A} \\
(3/2)(\alpha_s/\alpha_w) \approx 5.3 & \text{Scenario B}
\end{cases}
\end{equation}
In principle, precision measurements of weak-sector $\psi$-coupling
(e.g., through $\beta$-decay sensitivity to gravitational potential)
could discriminate between the two scenarios.  In practice, this
likely requires precision beyond current capabilities.

\paragraph{Theoretical preference.}
Scenario A (universal $\kappa_r = \alpha_r/4$) is more conservative
and relies on fewer assumptions (only the spectral action
normalization).  Scenario B introduces the additional hypothesis
that Berry-connection stiffness scales with the fundamental
representation dimension.  We adopt Scenario A as the baseline
and treat the additional factor of $n_r$ in Scenario B as a
theoretical upper bound.

%=============================================================================
\section{Implications for the Thorium-229 Nuclear Clock}
\label{sec:thorium}
%=============================================================================

The vacuum loading coefficient $\kappa_s$ affects the nuclear clock
predictions through the $\Lambda_{\rm QCD}$-dependent contributions
to the Th-229 isomer energy.

\paragraph{Unscreened estimate.}
Replacing $k_s$ with $\kappa_s$ in the nuclear clock enhancement:
\begin{equation}
\mathcal{R}_{\rm vac.load.}
= K_q \times A_\Lambda \times \frac{\kappa_s}{\kalpha}
= 10^4 \times 7.6 \times \frac{0.075}{8.5 \times 10^{-6}}
\approx 6.7 \times 10^8.
\end{equation}
This enormous unscreened value reinforces the conclusion that
terrestrial screening must reduce the effective coupling by many
orders of magnitude, consistent with the $\mu_{\rm LPI}$ screening
function approach.

\paragraph{Screened prediction.}
The screened nuclear clock enhancement is:
\begin{equation}
\mathcal{R}_{\rm screened} = \mathcal{R}_{\rm vac.load.}
\times \xi_{\rm LPI}^{\rm res}
\sim 6.7 \times 10^8 \times \xi_{\rm LPI}^{\rm res}.
\end{equation}
For $\xi_{\rm LPI}^{\rm res} \sim 10^{-6}$ (from clock bounds,
Sec.~\ref{sec:bounds}), $\mathcal{R}_{\rm screened} \sim 670$,
which is of the same order as the unscreened gauge-sector estimate
$\mathcal{R} \approx 170$ from the unified theory.  This provides
a self-consistency check.

%=============================================================================
\section{Summary and Conclusions}\label{sec:summary}
%=============================================================================

\begin{theorem}[Non-Abelian vacuum loading coefficient]
\label{thm:kappa-s}
Within the DFD gauge-emergence framework, the vacuum loading
coefficient for SU(3)$_c$ is:
\begin{equation}
\kappa_s = \frac{\alpha_s}{4} \qquad \text{(universal normalization)}
\end{equation}
or
\begin{equation}
\kappa_s = \frac{3\alpha_s}{4} \qquad
\text{(Berry-connection stiffness scaling)}
\end{equation}
with the universal normalization as the baseline.
\end{theorem}

\begin{tcolorbox}[colback=green!5!white, colframe=green!50!black,
title=Key Results]
\begin{enumerate}
\item \textbf{$\kappa_s = \alpha_s/4$} extends the U(1) result
$\kappa = \alpha/4$ to SU(3) using the same spectral action
normalization.

\item \textbf{At the hadronic scale} ($\alpha_s \approx 0.3$):
$\kappa_s = 0.075$, giving an effective coupling
$\kappa_s^{\rm eff} = \kappa_s \times A_\Lambda = 0.57$.

\item \textbf{Enhancement over $\kalpha$:}
$\kappa_s^{\rm eff}/\kalpha \approx 6.7 \times 10^4$, a factor
$\sim 260$ larger than the $k_s/\kalpha \approx 260$ from the
gauge--$\psi$ coupling alone.

\item \textbf{$\alpha$-tension gap:} With nuclear structure factors
($K_q \sim 10^4$), the total enhancement reaches
$\sim 6.7 \times 10^8$, more than sufficient to close the
$\alpha$-tension gap (required $\sim 10^7$) and demanding
screening with $\xi_{\rm LPI}^{\rm res} \lesssim 10^{-2}$.

\item \textbf{Experimental consistency:} The 2025 clock bound
$\dot{\Lambda}_{\rm QCD}/\Lambda_{\rm QCD} < 7 \times 10^{-17}\;
\mathrm{yr}^{-1}$ is satisfied for
$\xi_{\rm LPI}^{\rm res} \lesssim 10^{-6}$, consistent with
the screening framework.

\item \textbf{Discriminating test:} Precision weak-sector coupling
measurements could distinguish Scenario A ($\kappa_r = \alpha_r/4$)
from Scenario B ($\kappa_r = n_r\alpha_r/4$).
\end{enumerate}
\end{tcolorbox}

\paragraph{Status classification.}
\begin{center}
\renewcommand{\arraystretch}{1.2}
\begin{tabular}{lll}
\toprule
Result & Status & Basis \\
\midrule
$\kappa = \alpha/4$ (U(1)) & A- (near-rigorous) & Electroweak mixing (F118) \\
$\kappa_s = \alpha_s/4$ (Scenario A) & B+ (strong model) & Spectral action universality \\
$\kappa_s = 3\alpha_s/4$ (Scenario B) & B (model) & Berry-connection stiffness \\
Gap closure with screening & B (model) & Requires $\xi_{\rm LPI}^{\rm res}$ determination \\
\bottomrule
\end{tabular}
\end{center}

%=============================================================================
\section*{Appendix: Comparison of Enhancement Channels}
%=============================================================================

For reference, we collect all QCD enhancement mechanisms over the
baseline electromagnetic coupling $\kalpha$:

\begin{table}[h]
\centering
\caption{QCD enhancement channels.}\label{tab:channels}
\renewcommand{\arraystretch}{1.3}
\begin{tabular}{lccc}
\toprule
Channel & Formula & Value & Cumulative \\
\midrule
Baseline (EM) & $\kalpha = \alpha^2/(2\pi)$ & $8.5 \times 10^{-6}$ & 1 \\
Gauge hierarchy & $k_s/\kalpha = (\alpha_s/\alpha)^2$ & 260 & 260 \\
Vacuum loading & $\kappa_s/k_s = \pi/(2\alpha_s)$ & 13 (at $M_Z$) & 3,400 \\
Dim.\ transmutation & $A_\Lambda = 2\pi/(|b_3|\alpha_s)$ & 7.6 & 26,000 \\
Nuclear structure & $K_q$ & $10^4$ & $2.6 \times 10^8$ \\
\midrule
\multicolumn{3}{l}{\textbf{Required for $\alpha$-tension closure}}
& $\sim 10^7$ \\
\multicolumn{3}{l}{\textbf{Available (before screening)}}
& $\sim 10^{8\text{--}9}$ \\
\multicolumn{3}{l}{\textbf{Screening factor required}}
& $10^{-1}\text{--}10^{-2}$ \\
\bottomrule
\end{tabular}
\end{table}

\end{document}
